%% ECON 282E -- Homework 1, Part B
%% Reworked from PS2.lyx / PS2_IUBE.lyx (Quant Macro, HUST 2019).
%% Build:  python3 tools/build_docs.py homeworks-exams/HW1/hw1b.tex
%% With answers:  python3 tools/build_docs.py --solutions homeworks-exams/HW1/hw1b.tex
%% Answers live in solutions/ (gitignored); solution PDFs are gitignored too.
%% Neither must ever be committed.

\documentclass[11pt]{article}

\newif\ifsolutions \solutionsfalse

\usepackage[margin=1in]{geometry}
\usepackage{amsmath,amssymb,booktabs,enumitem,xcolor}
\usepackage[T1]{fontenc}
\usepackage{microtype}
\usepackage[colorlinks=true,linkcolor=blue!50!black,urlcolor=blue!50!black]{hyperref}

\newcommand{\solution}[1]{\ifsolutions\medskip
  {\color{blue!45!black}\textbf{Solution.} #1}\medskip\fi}

%% The answer key lives in solutions/ (gitignored, never published), so this
%% source carries no answers. --solutions pulls it in; a plain build ignores it.
\AtBeginDocument{\ifsolutions
  \IfFileExists{solutions/hw1b_key.tex}{\input{solutions/hw1b_key.tex}}
    {\GenericWarning{}{Answer key solutions/hw1b_key.tex not found}}\fi}

\title{\vspace{-2em}\textbf{ECON 282E --- Homework 1, Part B}\\[2mm]
       \large Continuous Choice, Stochastic Shocks, and Accuracy}
\author{Foundations of Macroeconomics \\ University of California, Riverside}
\date{Out: Thursday, October 15, 2026 \quad\textbullet\quad Due: Thursday, October 22, 2026}

\begin{document}
\maketitle
\thispagestyle{empty}

\noindent\rule{\linewidth}{0.4pt}\par
\small
\textbf{What to hand in.} One PDF, your code, and \texttt{AI\_LOG.md} --- together with Part A,
in a single submission. Question 3(d) feeds the anonymous accuracy leaderboard.

\textbf{Continuity with Part A.} Questions 2(d) and 2(e) extend the same deterministic model
and the same solver. Question 3 changes the calibration to a standard quarterly one and adds a
shock.
\normalsize\par
\noindent\rule{\linewidth}{0.4pt}\par

%=============================================================================
\section*{Question 2, continued. Getting off the grid (40 points)}

Same model and calibration as Part A: $u=\ln c$, $c+k'=Zk^{\alpha}$, $\delta=1$,
$\alpha=0.36$, $\beta=0.95$, $Z=1$.

\begin{enumerate}[label=(\alph*), leftmargin=*, start=4]

\item \textbf{(20 pts) Continuous choice.} Replace the grid search with a maximization over the
\emph{interval}, using golden section (or Brent) and evaluating the interpolated value function
off the grid. Keep $N=200$.
\begin{enumerate}[label=(\roman*)]
  \item Report the maximum relative policy error and the maximum unit-free Euler residual,
        evaluated at 500 capital levels that are \emph{not} grid points, against your Part A
        2(a) results.
  \item Report the number of objective evaluations and the wall time. How do they compare?
  \item Which interpolation scheme did you use for $\hat v$, and why does the choice matter for
        whether golden section is licensed at all?
\end{enumerate}

\solution{\keyBone}

\item \textbf{(20 pts) A second road.} Solve the same model again by \emph{either} time
iteration on the Euler equation \emph{or} the endogenous grid method. Then produce one table
comparing all four solvers --- Part A 2(a), Part A 2(c), and 2(d) and 2(e) here --- on
iterations, wall time, maximum relative policy error and maximum off-grid Euler residual.
Which solver would you use in practice, and for what reason?

\solution{\keyBtwo}

\end{enumerate}

%=============================================================================
\section*{Question 3. A stochastic, quarterly model (60 points)}

Now the standard quarterly calibration:
\[
  \alpha=0.33,\qquad \beta=0.99,\qquad \delta=0.025,\qquad
  \ln z'=\rho\ln z+\sigma\varepsilon',\quad \rho=0.95,\ \sigma=0.007 ,
\]
with $c+k'=zk^{\alpha}+(1-\delta)k$ and $u=\ln c$. There is \textbf{no closed-form policy} for
this model --- only a closed-form steady state,
$k^{*}=\big(\alpha\beta/(1-\beta(1-\delta))\big)^{1/(1-\alpha)}$.

\begin{enumerate}[label=(\alph*), leftmargin=*]

\item \textbf{(15 pts) Where $\Pi$ comes from.} Implement \emph{both} Tauchen (1986) and
Rouwenhorst (1995). For $n\in\{5,9,15\}$, report each chain's implied persistence and
unconditional standard deviation against the targets $\rho$ and $\sigma/\sqrt{1-\rho^{2}}$.
Then repeat the table at $\rho=0.99$. Which method would you use, and why?

\solution{\keyBthree}

\item \textbf{(15 pts) Solve it.} Using the Rouwenhorst chain with $n=7$, solve the model with
your best solver from Question 2. Report the grid, tolerance, sweeps and wall time, and verify
that your policy crosses the 45-degree line near the analytical $k^{*}$ in the middle
productivity state. How close is the crossing? If it is further off than one grid spacing,
diagnose why --- the answer is not that your solver is broken.

\solution{\keyBfour}

\item \textbf{(15 pts) Simulate and match moments.} Simulate 2,000 quarters (discard the first
500), HP-filter $\ln Y$, $\ln C$ and $\ln I$ with $\lambda=1600$, and report
$\sigma(Y)$, $\sigma(C)/\sigma(Y)$, $\sigma(I)/\sigma(Y)$, $\mathrm{corr}(C,Y)$,
$\mathrm{corr}(I,Y)$ and the first-order autocorrelation of $Y$. Compare with the US
post-war targets $\sigma(Y)\approx1.5$--$2.0\%$, $\sigma(C)/\sigma(Y)\approx0.5$--$0.8$,
$\sigma(I)/\sigma(Y)\approx3$--$4$, both correlations $\approx0.8$--$0.9$.
\emph{An HP filter is provided in Lab 3, Part 7; you may use it or your own.}
Then re-run with the \emph{Tauchen} chain and report how much the moments move.

\solution{\keyBfive}

\item \textbf{(15 pts) Accuracy, and the leaderboard.} Compute unit-free Euler residuals
$\mathcal{E}=|1-\tilde c/c|$ at 1,000 capital levels drawn uniformly, \emph{off} the solution
grid, across all productivity states. Report $\log_{10}$ of the mean and the maximum, and plot
the residual against $k$. Where are the errors largest, and why?

Submit your maximum $\log_{10}\mathcal{E}$, together with your grid size, tolerance and wall
time, for the anonymous class leaderboard.

\solution{\keyBsix}

\end{enumerate}

%=============================================================================
\section*{Bonus (up to 10 points). Endogenous labour}

Add a labour choice with $u(c,n)=\ln c-\psi\frac{n^{1+1/\nu}}{1+1/\nu}$ and
$y=zk^{\alpha}n^{1-\alpha}$. Calibrate $\psi$ so that average hours in the simulation are
$\bar n=0.3$. Report how $\sigma(Y)$ and $\sigma(I)/\sigma(Y)$ change relative to Question 3,
and say which of the target moments the model now gets closer to.

\solution{\keyBseven}

\vfill
\noindent\rule{\linewidth}{0.4pt}\par
\small
\textbf{Reporting standard.} With every result: the method, the grid (points, spacing, bounds),
the chain and its \emph{implied} $\rho$ and $\sigma_z$, the tolerance, the dtype, maximum and
mean Euler residuals on points not used to solve, and the wall time with the hardware. A number
without these is not yet a result.
\normalsize

\end{document}
